% =============================================================================
% 037 / 068
% Status: raw
% =============================================================================
% Notes:
%
% Source question:
% 你在書面語那篇文中提到：「德語250年......外加網際網路的傳播可以讓新吳語在20年左右普及......」「吳語如果沒有新的發展就談不上維持和保護。未來，無法做到言文一致的語言將很快在AI時代消亡。」
% 
% 閣下眼中，是不是在新時代與AI當中，
% 看到了復興吳語 建構書面吳語新的可能性？
% =============================================================================

\chapter[你在書面語那篇文中提到：「德語250年......外加網際網路的傳播可以讓新吳語在20年左右普及......」「吳語如果沒有新的發展就談不上維持和保護。未來，無法做到言文一致的語言將很快在AI時代消亡。」 閣下眼中，是不是在新時代與AI當中， 看到了復興吳語 建構書面吳語新的可能性？]{你在書面語那篇文中提到：「德語250年......外加網際網路的傳播可以讓新吳語在20年左右普及......」「吳語如果沒有新的發展就談不上維持和保護。未來，無法做到言文一致的語言將很快在AI時代消亡。」 \\[0.35em] 閣下眼中，是不是在新時代與AI當中， \\[0.35em] 看到了復興吳語 建構書面吳語新的可能性？}
\qaanswer{
吳語不發展，不進行語法整理和詞彙疏理、構建，不透過書面語定型，不出30年就會徹底滅絕。所謂的8000多萬吳語使用者只是一個30年前統計的虛數。現下除去已經逝去的高齡吳語使用者，和現在絕大多數中老年藍青話使用者外，真正還會原生吳語，掌握一定量吳語詞彙的人估計100萬都不會有了。AI時代只是在整理和構建吳語時能夠用到的更為高效、便捷的工具，語言構建這件事本身如果沒人做的話到2030年代這個時候，區別於漢語官話的原生吳語使用者可能就已經不存在了。
}

